\documentclass[10pt,a4paper]{amsart}
\usepackage[margin=19mm]{geometry}
\usepackage{amsmath,amssymb,mathtools}
\usepackage[scr=rsfs,cal=euler]{mathalfa}
\usepackage{xfrac}
\usepackage[unicode,hidelinks,bookmarks=false]{hyperref}
\newcommand{\ie}{i.e.\ }
\newcommand{\cf}{cf.\ }
\newcommand{\resp}{respectively\ }
\newcommand{\Subsection}{\subsection}
\providecommand{\nobreakdash}{\nobreak\hbox{-}}
\setlength{\emergencystretch}{2em}
\multlinegap0pt
\usepackage{amssymb}
\usepackage{xcolor}
\usepackage{algorithm}\usepackage{algpseudocode}
\usepackage[shortlabels]{enumitem}
\usepackage{float}
\newtheorem{lemma}{Lemma}
\title{Optimal Harvesting of Size-Structured Populations with Environmental Feedback and Fixed Recruitment Flux}
\author{Louis Shuo Wang, Jiguang Yu}
\date{Source: arXiv:2607.03790v1 (CC BY 4.0)}
\begin{document}
\maketitle

\noindent\emph{Source and licence.} Optimal Harvesting of Size-Structured Populations with Environmental Feedback and Fixed Recruitment Flux. Version arXiv:2607.03790v1 (CC BY 4.0), distributed by arXiv under the Creative Commons Attribution 4.0 International licence, http://creativecommons.org/licenses/by/4.0/, as recorded in the arXiv licence metadata for this exact version and shown on the abstract page https://arxiv.org/abs/2607.03790v1. Full licence notice: \texttt{licenses/arxiv-2607.03790-CC-BY-4.0.txt}.\par\smallskip

\noindent\textit{Editorial scope.} Counted unit: the article's lemma lem:app\_bv\_composition (source0.tex lines 1514--1539). The article's complete original proof of it is delivered with it (source0.tex lines 1541--1628, one \texttt{proof} environment of 88 lines). No mathematics is written, repaired or re-derived: the statement and the proof are the article's own lines, in the article's own notation, and no retained line is substituted or re-flowed.  No further source-local context is needed: every cross-reference of every retained line resolves inside the two delivered spans.  Nothing was omitted from any retained span, so the package carries no omission record.  No retained line contains a citation, so the wrapper carries no bibliography and no bibliography entry.  No retained line contains a figure environment, an image-inclusion command or drawing code, so no image is delivered and the package declares no asset dependency.  Editorial formatting only, in addition: an AMS \texttt{amsart} preamble that restates the article's own declarations and macros used by the retained lines, namely the environments \texttt{lemma} on separate counters, together with the standard packages those lines need.  The retained environments print in this document as Lemma 1 in order of appearance, whereas the article prints the same items with the numbering of its own full document.  The complete native source of the article as distributed in the arXiv e-print for this exact version is bundled unchanged as \texttt{sources/204421/source0.tex}.
\begin{lemma}[\(BV\)-composition estimate]
\label{lem:app_bv_composition}
Let \(w\in BV(\Omega)\), and extend it to
\(\overline w\in BV(\mathbb R)\) by constant continuation,
so that
\[
\operatorname{TV}_{\mathbb R}(\overline w)
=
\operatorname{TV}_{\Omega}(w).
\]
Let
\(\phi_1,\phi_2:\Omega\to\Omega\)
be monotone bi-Lipschitz maps satisfying
$|\phi_i'|\ge J_m>0$ a.e., and
$\|\phi_1-\phi_2\|_{L^\infty(\Omega)}
\le d$.
Then
\[
\int_\Omega
|w(\phi_1(l))-w(\phi_2(l))|
\,dl
\le
\frac{d}{J_m}
\operatorname{TV}_\Omega(w).
\]
\end{lemma}

\begin{proof}
By the change of variables
\(\xi=\phi_1(l)\),
the lower bound on \(\phi_1'\) yields
$dl \le J_m^{-1}\,d\xi$,
hence
\[
\int_\Omega
|w(\phi_1(l))-w(\phi_2(l))|
\,dl
\le
\frac1{J_m}
\int_{\phi_1(\Omega)}
|\overline w(\xi)-\overline w(\psi(\xi))|
\,d\xi,
\]
where $\psi(\xi) := \phi_2(\phi_1^{-1}(\xi))$.
Since $|\psi(\xi)-\xi| \le d$,
it remains to estimate the last integral.
For the precise representative of a one-dimensional
\(BV\) function,
the distributional derivative
\(D\overline w\)
is a finite Radon measure and
\[
\overline w(b)-\overline w(a)
=
D\overline w((a,b])
\]
for every \(a<b\)
(see Ambrosio--Fusco--Pallara).
Therefore,
\[
|\overline w(\xi)-\overline w(\psi(\xi))|
\le
|D\overline w|
\!\left(
[\min\{\xi,\psi(\xi)\},
\max\{\xi,\psi(\xi)\}]
\right).
\]
Integrating over \(\phi_1(\Omega)\) and applying Fubini,
\[
\int_{\phi_1(\Omega)}
|\overline w(\xi)-\overline w(\psi(\xi))|
\,d\xi \le
\int_{\mathbb R}
\left(
\int_{\phi_1(\Omega)}
\mathbf 1_{
[\min(\xi,\psi(\xi)),
\max(\xi,\psi(\xi))]
}(y)
\,d\xi
\right)
\,d|D\overline w|(y).
\]
Since every interval
joining
\(\xi\)
and
\(\psi(\xi)\)
has length at most \(d\),
\[
\int_{\phi_1(\Omega)}
\mathbf 1_{
[\min(\xi,\psi(\xi)),
\max(\xi,\psi(\xi))]
}(y)
\,d\xi
\le
d,
\]
for every \(y\in\mathbb R\).
Hence
\[
\int_{\phi_1(\Omega)}
|\overline w(\xi)-\overline w(\psi(\xi))|
\,d\xi
\le
d\,|D\overline w|(\mathbb R)
=
d\,\operatorname{TV}_{\mathbb R}(\overline w).
\]
Since $\operatorname{TV}_{\mathbb R}(\overline w)
=
\operatorname{TV}_\Omega(w)$, the desired estimate follows.
\end{proof}

\end{document}
